% Lösungen Einheit 2
\einheitenkopf*{Einheit 2 von 3 · Kreisfläche}
\erg{11}{a) $A = \pi \cdot r^2$ \quad b) $u = \pi \cdot d$ \quad c) $A = \pi \cdot r^2$ (erst halbieren) \quad d) $u = \pi \cdot d$ \quad e) $A = \pi \cdot r^2$}
\erg{12}{a) 3,1\,cm² \quad b) 153,9\,cm² \quad c) 254,5\,m² \quad d) 452,4\,cm² \quad e) $r = 8$\,cm, $A \approx 201{,}1$\,cm² \quad f) 19,6\,cm²}
\erg{13}{a) $\frac12 \cdot \pi \cdot 8^2 \approx 100{,}5$\,cm² \quad b) $\frac14 \cdot \pi \cdot 11^2 \approx 95{,}0$\,cm² \quad c) $r = 4{,}5$\,m, $A \approx 31{,}8$\,m² \quad d) $A = \frac14 \cdot \pi \cdot r^2$ \quad e) ein Dreiviertelkreis (drei Viertel des Kreises)}
\erg{14}{\begin{minipage}[t]{0.9\linewidth}\vspace{0pt}\sachtabelle{lcccc}{ & Radius $r$ & Durchmesser $d$ & Umfang $u$ & Flächeninhalt $A$}{a) & 1,5\,m & 3\,m & 9,4\,m & 7,1\,m² \\ b) & 15\,cm & 30\,cm & 94,2\,cm & 706,9\,cm² \\ c) & 6,4\,cm & 12,7\,cm & 40\,cm & 127,3\,cm²}\end{minipage}}
\erg{15}{a) 3,1\,cm \quad b) 8,0\,m \quad c) $r = \sqrt{1000 : \pi} \approx 17{,}84$\,cm, $d \approx 35{,}7$\,cm \quad d) $r = 3{,}66 : 2 = 1{,}83$\,m; $A = \pi \cdot 1{,}83^2 \approx 10{,}52$\,m²}
\erg{16}{a) Mia rechnet $r^2$ als $2 \cdot r$. Richtig: $A = \pi \cdot 17^2 = \pi \cdot 289 \approx 907{,}9$\,cm² \quad b) $0{,}8^2 = 0{,}64$; $A \approx 2{,}0$\,m²}
\erg{17}{a) Flächeninhalt: $\pi \cdot 25^2 \approx 1963{,}5$\,cm² \quad b) Umfang: $2 \cdot \pi \cdot 25 \approx 157{,}1$\,cm \quad c) Umfang: $\pi \cdot 4{,}2 \approx 13{,}2$\,m \quad d) Flächeninhalt: $r = 2{,}1$\,m, $A \approx 13{,}9$\,m²}
\erg{18}{a) etwa der halbe Umfang, also $\pi \cdot r$ (die Bögen der Stücke liegen je zur Hälfte unten und oben) \quad b) etwa $r$ \quad c) Die Figur ist fast ein Rechteck: $A \approx \pi \cdot r \cdot r = \pi \cdot r^2$ \quad d) Nein: $(2r)^2 = 4r^2$, die Fläche wird viermal so groß.}
\erg{19}{a) klein: $r = 13$\,cm, $A \approx 530{,}9$\,cm²; groß: $r = 18$\,cm, $A \approx 1017{,}9$\,cm² \quad b) klein: $8{,}50\,€ : 5{,}309 \approx 1{,}60$\,€; groß: $14{,}00\,€ : 10{,}179 \approx 1{,}38$\,€ \quad c) Die große Pizza ist günstiger: 100\,cm² kosten dort nur etwa 1,38\,€ statt 1,60\,€.}
